\documentclass{article}
\usepackage[UTF8]{ctex}
\usepackage{amsmath,amssymb,amsfonts}
\usepackage{enumitem}
\usepackage{geometry}
\geometry{a4paper, margin=1in}
% === 错题本排版优化（不参与题目标识） ===
\setlength{\parindent}{0pt}
\setlength{\parskip}{0.35em plus 0.10em minus 0.08em}
\setlength{\emergencystretch}{3em}
\setlist{itemsep=0.25em, topsep=0.40em, parsep=0pt}
\allowdisplaybreaks[2]
% === 错题本排版优化结束 ===
\begin{document}
\section*{南昌大学2025年高等代数真题}
\section*{112．南昌大学2025年高等代数试卷}

\subsection{题目1}
1．（15 分）设 \(y \neq z\) ，计算 \(n\) 阶行列式

\[
D_{n}=\left|\begin{array}{ccccc}
x_{1} & y & \cdots & y & y \\
z & x_{2} & \cdots & y & y \\
\vdots & \vdots & & \vdots & \vdots \\
z & z & \cdots & x_{n-1} & y \\
z & z & \cdots & z & x_{n}
\end{array}\right| .
\]

\subsection{题目2}
2．（15 分）已知 \(n\) 为奇数，证明：\((x+y)(y+z)(x+z)\) 整除 \((x+y+z)^{n}-x^{n}-y^{n}-z^{n}\) ．

\subsection{题目3}
3．（15 分）已知方程组\\
（I）\(\left\{\begin{array}{l}7 x_{1}-6 x_{2}+3 x_{3}=b ; \\ 8 x_{1}-9 x_{2}+a x_{4}=7 .\end{array}\right.\)\\
方程组（II）的通解为

\[
(1,1,0,0)^{\prime}+t_{1}(1,0,-1,0)+t_{2}(2,3,0,1)^{\prime} .
\]

若（I）与（II）有无穷多公共解，求 \(a, b\) 的值与公共解．

\subsection{题目4}
4．（15 分）已知 \(A\) 是 \(n\) 阶幂零阵，\(B\) 是 \(n\) 阶矩阵，\(A B+B A=B\) ，证明：\(B=O\) 。

\subsection{题目5}
5．（15 分）求 \(n\) 阶矩阵 \(A=\left(\begin{array}{cccc}a & a & \cdots & a \\ 0 & a & \cdots & a \\ \vdots & \vdots & \ddots & \vdots \\ 0 & 0 & \cdots & a\end{array}\right)\) 的若尔当标准形．

\subsection{题目6}
6．（15 分）设 \(V\) 是实数域上连续函数构成的实线性空间，证明： \(1, \cos x, \cos 2 x, \cdots, \cos n x\) 线性无关．

\subsection{题目7}
7．（15 分）\(n\) 阶矩阵 \(A\) 的全部特征值为 \(\lambda_{1}, \lambda_{2}, \cdots, \lambda_{n}\) ，求 \(2 n\) 阶矩阵

\[
\left(\begin{array}{cc}
A & A^{2} \\
A^{2} & A
\end{array}\right)
\]

的全部特征值。

\subsection{题目8}
8．（15 分）\(A, B\) 都为 \(n\) 阶矩阵，且 \(A B=B A=O, r(A)=r\left(A^{2}\right)\) ．证明：\(r(A+B)=r(A)+r(B)\) ．

\subsection{题目9}
9．（15 分）\(J=J_{n}\left(\lambda_{0}\right)\) 是特征值为 \(\lambda_{0}\) 的 \(n\) 阶若尔当块，证明：和 \(J\) 可乘法交换的 \(n\) 阶矩阵必可表示为 \(J\) 次数不超过 \(n-1\) 的多项式。

\subsection{题目10}
10．（15 分）设 \(f(x)=a_{m} x^{m}+a_{m-1} x^{m-1}+\cdots+a_{1} x+a_{0}\) 以 \(n\) 阶矩阵 \(A\) 为根，\(a_{m}>\sum_{i=1}^{m-1}\left|a_{i}\right|\) ，证明：矩阵方程 \(2 X+A X=X A^{2}\) 只有零解．
\end{document}
